\documentclass[
12pt,
openright,
oneside,
a4paper,
brazil
]{abntex2}

% ------------------------------------------------
% PACOTES
% ------------------------------------------------

\usepackage[utf8]{inputenc}
\usepackage[T1]{fontenc}
\usepackage{lmodern}
\usepackage{indentfirst}
\usepackage{graphicx}
\usepackage{microtype}
\usepackage{setspace}
\usepackage{longtable}
\usepackage{float}

\usepackage{hyperref}

% ------------------------------------------------
% DADOS DO TRABALHO
% ------------------------------------------------

\titulo{Sistema de Tutoria Baseado em IA Educativa para Avaliação Formativa em Ciências Humanas}

\autor{
Luis Henrique Gomes Zortea \\
Marcos Vinicius Silva Torres
}

\local{Vila Velha -- ES}
\data{2026}

\instituicao{
Universidade Vila Velha \\
Departamento de Computação
}

\tipotrabalho{Artigo}
\preambulo{
Artigo apresentado ao Curso de Ciência da Computação da Universidade Vila Velha como parte dos requisitos para obtenção do grau de Bacharel em Ciência da Computação.
}

\orientador{Prof.ª Fabiana Bravim de Freitas}

% ------------------------------------------------
% INÍCIO DO DOCUMENTO
% ------------------------------------------------

\begin{document}

% CAPA
\imprimircapa

% FOLHA DE ROSTO
\imprimirfolhaderosto*


% ------------------------------------------------
% RESUMO
% ------------------------------------------------

\begin{resumo}

Este trabalho apresenta o desenvolvimento da plataforma NextEducation, um sistema web de avaliação formativa voltado ao apoio ao estudo de Ciências Humanas para o Exame Nacional do Ensino Médio (ENEM). A plataforma integra questões educacionais, análise de desempenho e geração automática de feedback didático por meio de modelos de linguagem de grande escala acessados via API Groq. A aplicação foi desenvolvida com o framework Next.js e TypeScript, utilizando Tailwind CSS para a interface e MongoDB para persistência de dados. A proposta busca transformar a correção objetiva de questões em uma experiência de tutoria inteligente, oferecendo explicações contextualizadas para os erros dos estudantes e auxiliando na identificação de dificuldades conceituais. Como resultado, foi desenvolvido um protótipo funcional composto por módulos de autenticação, atividades, mini provas, dashboard de desempenho e tutor IA. Os testes técnicos preliminares indicam viabilidade da integração com a API Groq, embora ainda sejam necessárias avaliações pedagógicas sistemáticas com estudantes e especialistas para verificar a qualidade, a clareza e a confiabilidade dos feedbacks gerados.

\textbf{Palavras-chave}: Inteligência Artificial. ENEM. Ciências Humanas. Educação a Distância. Next.js. MongoDB.

\end{resumo}

% ------------------------------------------------
% ABSTRACT
% ------------------------------------------------

\begin{resumo}[Abstract]

This work presents the development of NextEducation, a web-based formative assessment platform designed to support Human Sciences students preparing for the Brazilian ENEM exam. The platform integrates educational questions, performance analysis, and automatic pedagogical feedback generated by large language models through the Groq API. The system was developed using Next.js and TypeScript, with Tailwind CSS for the interface and MongoDB for data persistence. The proposal aims to transform traditional answer checking into an intelligent tutoring experience by providing contextualized explanations for students’ mistakes and supporting the identification of conceptual difficulties. The resulting prototype includes authentication, learning activities, mini exams, a performance dashboard, and an AI tutor. Preliminary technical tests indicate the feasibility of integrating the platform with the Groq API, although systematic pedagogical evaluations with students and experts are still needed to assess the quality, clarity, and reliability of the generated feedback.

\textbf{Keywords}: Artificial Intelligence. ENEM. Human Sciences. Distance Education. Next.js. MongoDB.

\end{resumo}

% ------------------------------------------------
% LISTAS E SUMÁRIO
% ------------------------------------------------

\pdfbookmark[0]{\contentsname}{toc}
\tableofcontents*
\cleardoublepage

% ------------------------------------------------
% TEXTO PRINCIPAL
% ------------------------------------------------

\textual

% ------------------------------------------------
% INTRODUÇÃO
% ------------------------------------------------

  \chapter{Introdução}

O cenário educacional contemporâneo tem sido transformado pela integração de sistemas computacionais que visam otimizar a aquisição de conhecimento e o acompanhamento pedagógico. No entanto, um dos maiores desafios enfrentados por estudantes em exames de larga escala, como o Exame Nacional do Ensino Médio (ENEM), é a retenção de conteúdos a longo prazo e o entendimento profundo das áreas do conhecimento. Essa retenção é frequentemente prejudicada por métodos de estudo massificados e lineares, que priorizam a memorização temporária em detrimento do aprendizado significativo. A ciência cognitiva aponta que a Repetição Espaçada (Spaced Repetition) \cite{ebbinghaus, cepeda, kang} é uma das técnicas mais eficazes para mitigar a curva do esquecimento, propondo revisões em intervalos matematicamente calculados para consolidar a memória.

Apesar da eficácia inegável dessa técnica, a sua aplicação prática no estudo autônomo exige ferramentas não apenas para gerenciar o agendamento de revisões, mas principalmente para oferecer um suporte qualitativo ao estudante durante a resolução de problemas. É neste ponto que a avaliação formativa demonstra sua relevância. Conforme apontado por Black e Wiliam \cite{blackwilliam} e Sadler \cite{sadler}, a avaliação formativa transcende a simples medição de desempenho, caracterizando-se por envolver feedback contínuo, identificação precisa de dificuldades e fornecimento de orientação estruturada para a melhoria constante da aprendizagem.

Nesse contexto, este projeto propõe o desenvolvimento de uma plataforma web inovadora que utiliza Inteligência Artificial Generativa para transformar a simples conferência de gabaritos em uma experiência robusta de tutoria inteligente \cite{kasneci2023chatgpt}. Ao integrar os princípios da avaliação formativa com os avanços dos modelos de linguagem de grande escala (LLMs), a ferramenta não apenas indica o erro do aluno, mas fornece uma explicação contextualizada que atua como mediação pedagógica em tempo real, favorecendo a correção de rotas cognitivas e a consolidação do conhecimento em Ciências Humanas.

\section{Justificativa}

A justificativa deste trabalho reside na carência de soluções open source que integrem algoritmos de repetição espaçada, notificações em aplicativos de mensagens e o poder de processamento de linguagem natural dos modelos de linguagem modernos. Enquanto plataformas tradicionais focam em cálculos exatos, o estudo de Humanas exige interpretação e contextualização, áreas nas quais modelos de linguagem podem auxiliar na geração de explicações contextualizadas e feedbacks automatizados para apoio ao aprendizado.

Ao integrar notificações proativas em mensageiros e utilizar o dataset ENEM-Benchmark \cite{maritaca}, o sistema busca aumentar o engajamento do aluno com as atividades de estudo e favorecer a continuidade do aprendizado, contribuindo para o desenvolvimento da autonomia e oferecendo uma infraestrutura tecnológica robusta para o aprendizado autodirigido.

\section{Objetivos}

\subsection{Objetivo Geral}

O objetivo geral deste trabalho é projetar, implementar e avaliar preliminarmente uma plataforma web de avaliação formativa para Ciências Humanas no contexto do ENEM, integrando questões educacionais, análise de desempenho e feedback didático gerado por IA generativa, com o propósito de apoiar o estudo autodirigido e a identificação de dificuldades conceituais dos estudantes.

A pesquisa foi orientada pelas seguintes questões:

\begin{itemize}
\item QP1: Como integrar uma base de questões do ENEM, análise de desempenho e modelos de linguagem em uma plataforma web de avaliação formativa?

\item QP2: Em que medida um modelo de linguagem de grande escala é capaz de gerar feedbacks didáticos coerentes com o gabarito oficial e pedagogicamente úteis para estudantes de Ciências Humanas?

\item QP3: Quais limitações técnicas e pedagógicas emergem no uso de IA generativa para correção, justificativa e tutoria em questões objetivas do ENEM?
\end{itemize}

\subsection{Objetivos Específicos}

Os objetivos específicos deste trabalho são:

\begin{itemize}
\item Levantar e especificar os requisitos funcionais e não funcionais da plataforma;
\item Selecionar, tratar e organizar questões de Ciências Humanas provenientes do dataset ENEM-Benchmark;
\item Desenvolver uma plataforma web com módulos de autenticação, atividades, mini provas, dashboard de desempenho e tutor IA;
\item Integrar a plataforma a um modelo de linguagem por meio da API Groq;
\item Gerar feedbacks didáticos automatizados a partir da resposta do estudante e do gabarito oficial;
\item Avaliar preliminarmente a estabilidade técnica do protótipo e a qualidade pedagógica dos feedbacks gerados.
\end{itemize}

% ------------------------------------------------
% TRABALHOS RELACIONADOS
% ------------------------------------------------

\chapter{Trabalhos Relacionados}

\section{Projetos Educacionais Relacionados}

Nesta seção, são apresentados projetos e estudos que fundamentam o estado da arte em sistemas de aprendizagem espaçada e tutoria inteligente, servindo de base comparativa para a arquitetura proposta neste trabalho.

\subsection{Projeto 1}
O projeto LECTOR, proposto por Zhao\cite{zhao2025lector} (2025), apresenta um algoritmo que integra Modelos de Linguagem de Larga Escala (LLMs) à técnica de repetição espaçada para realizar análises semânticas de conteúdo. O diferencial deste sistema reside na sua capacidade de adaptação dinâmica em exames de alta performance, utilizando a IA para ajustar os ciclos de revisão com base na compreensão profunda do estudante, e não apenas em acertos binários.

\subsection{Projeto 2}
No cenário nacional, Damasceno\cite{damasceno2021stuart} (2021) desenvolveu o STUART (Sistema de Tutoria Inteligente Artificial), focado em aumentar a escalabilidade no ensino a distância. O STUART utiliza uma arquitetura de tutoria inteligente integrada a chatbots para fornecer interações didáticas e recomendações de conteúdo aos alunos. Este trabalho corrobora a viabilidade do uso de mensageiros como interface de suporte pedagógico, embora o presente TCC avance ao aplicar essa lógica especificamente ao contexto do ENEM e das Ciências Humanas via IA generativa.

\subsection{Projeto 3}
Complementando a visão de sistemas de código aberto, Seabra\cite{seabra2023memorizapp} (2023) apresenta o MemorizApp, uma plataforma \textit{open-source} de repetição espaçada baseada em \textit{flashcards} digitais. O projeto foca em superar limitações de usabilidade de ferramentas tradicionais (como o Anki), propondo uma interface simplificada para o aprendizado baseado em cartões. Enquanto o MemorizApp foca na gestão eficiente da repetição, a solução proposta neste artigo diferencia-se pela camada de inteligência que fornece explicações automáticas para os erros dos usuários.


% ------------------------------------------------
% Desenvolvimento Temático
% ------------------------------------------------

\chapter{Metodologia}

Este trabalho caracteriza-se como uma pesquisa aplicada, de natureza exploratória, orientada ao desenvolvimento e avaliação preliminar de um artefato computacional educacional. O artefato desenvolvido corresponde a uma plataforma web de avaliação formativa para Ciências Humanas, integrada a modelos de linguagem por meio da API Groq \cite{groq}.

A condução da pesquisa foi metodologicamente organizada em cinco etapas distintas, detalhadas a seguir:

\textbf{Etapa 1: Levantamento de Requisitos.} Esta fase consistiu na identificação e documentação das necessidades técnicas e pedagógicas da plataforma, com foco em estudantes do ENEM. Os requisitos foram divididos em funcionais (o que o sistema deve fazer) e não funcionais (como o sistema deve se comportar), conforme ilustrado na Tabela \ref{tab:requisitos}.

\begin{table}[H]
    \centering
    \caption{Requisitos Funcionais e Não Funcionais do Sistema}
    \begin{tabular}{|p{0.45\textwidth}|p{0.45\textwidth}|}
        \hline
        \textbf{Requisitos Funcionais} & \textbf{Requisitos Não Funcionais} \\
        \hline
        - Cadastro e autenticação de usuários & - Interface responsiva e acessível \\
        - Resolução de questões de simulado & - Baixa latência na resposta do Tutor IA \\
        - Geração de feedback em tempo real & - Arquitetura escalável baseada na nuvem \\
        - Acompanhamento de métricas no Dashboard & - Persistência segura em banco NoSQL \\
        \hline
    \end{tabular}
    \label{tab:requisitos}
\end{table}

\textbf{Etapa 2: Curadoria de Dados.} Envolveu a seleção, tratamento e estruturação de questões provenientes do dataset oficial ENEM-Benchmark \cite{maritaca}. O foco foi estrito em itens da área de Ciências Humanas. Os dados, disponibilizados no formato JSONL, foram inseridos em um banco de dados MongoDB para acesso rápido via aplicação.

\textbf{Etapa 3: Desenvolvimento da Plataforma.} Correspondeu à codificação e estruturação arquitetônica da solução. Utilizou-se o framework Next.js (App Router), viabilizando os módulos de autenticação, atividades, mini provas, painel de desempenho (dashboard) e interface do tutor IA. Optou-se pelo Tailwind CSS para garantir consistência visual. A Figura \ref{fig:arquitetura} ilustra o desenho de alto nível da arquitetura da solução, cujos termos principais equivalem a: \textit{Client/Frontend} (Cliente/Frontend); \textit{Server/Backend} (Servidor/Backend); \textit{External Database} (Banco de Dados Externo); \textit{External AI Services} (Serviços Externos de IA); \textit{HTTP Requests} (Requisições HTTP); e \textit{AI Inference} (Inferência por IA). 

A arquitetura baseia-se em um modelo cliente-servidor, onde a camada de Cliente é responsável pela interface e interação direta com o aluno. Esta camada envia Requisições HTTP para a camada de Servidor, que concentra a lógica de negócio e as \textit{Server Actions}. O servidor atua como intermediário seguro: ele se conecta ao Banco de Dados Externo (MongoDB) para gravar o progresso e recuperar questões, e simultaneamente se comunica com os Serviços Externos de IA (Groq API) para solicitar a Inferência por IA, recebendo e formatando o feedback didático antes de enviá-lo de volta ao cliente.

\begin{figure}[H]
    \centering
    \includegraphics[width=0.9\textwidth]{imagens/arquitetura_nexteducation_pt.png}
    \caption{Desenho arquitetural de alto nível da solução NextEducation.}
    \label{fig:arquitetura}
\end{figure}

\textbf{Etapa 4: Integração com Inteligência Artificial.} Focou-se na comunicação backend entre a plataforma e o modelo de linguagem Llama 3.3, provido pela API da Groq. O processo exigiu o desenvolvimento de \textit{prompts} pedagógicos especializados para orientar a IA a justificar os erros e contextualizar os acertos.

\textbf{Etapa 5: Avaliação Técnica e Pedagógica Preliminar.} Consistiu no teste do protótipo desenvolvido. Para cada uma das 20 questões selecionadas da área de Ciências Humanas, foi simulada uma resposta de estudante e gerado um feedback didático pelo Tutor IA. Dessa forma, o \textit{corpus} de avaliação foi composto por 20 feedbacks, sendo um feedback correspondente a cada questão. Os feedbacks foram avaliados consensualmente pelos próprios autores, utilizando uma escala Likert de 1 a 5, em que 1 representa avaliação muito inadequada e 5 representa avaliação muito adequada. Foram considerados quatro critérios: correção conceitual, aderência ao gabarito oficial, clareza didática e contextualização em Ciências Humanas.

A fim de tornar a avaliação objetiva e reprodutível, foi estabelecida uma rubrica de avaliação detalhada, que define os parâmetros qualitativos para cada nível da escala, conforme apresentado na Tabela \ref{tab:rubrica}.

\begin{table}[H]
    \centering
    \caption{Rubrica de Avaliação Pedagógica dos Feedbacks Gerados pela IA}
    \begin{tabular}{|p{0.05\textwidth}|p{0.85\textwidth}|}
        \hline
        \textbf{Nota} & \textbf{Descrição do Nível de Adequação} \\
        \hline
        1 & \textbf{Muito inadequado:} O feedback apresenta erro conceitual grave, não se relaciona adequadamente com a questão, contradiz o gabarito oficial ou pode induzir o estudante ao erro. \\
        \hline
        2 & \textbf{Inadequado:} O feedback apresenta limitações importantes, como explicação incompleta, linguagem excessivamente técnica ou baixa aderência ao gabarito oficial. \\
        \hline
        3 & \textbf{Parcialmente adequado:} O feedback é parcialmente correto, mas possui lacunas, generalizações, pouca profundidade didática ou contextualização superficial. \\
        \hline
        4 & \textbf{Adequado:} O feedback é correto, compreensível e útil para o estudante, reconhece o gabarito, mas ainda pode apresentar pontos de melhoria na clareza. \\
        \hline
        5 & \textbf{Muito adequado:} O feedback é conceitualmente correto, perfeitamente claro, bem estruturado, aderente ao gabarito e altamente contextualizado em Ciências Humanas. \\
        \hline
    \end{tabular}
    \label{tab:rubrica}
\end{table}

O critério de aprovação preliminar para o feedback de uma questão foi obter pontuação média igual ou superior a 4 em todos os critérios listados.

% ------------------------------------------------
% Tecnologias
% ------------------------------------------------

\chapter{Desenvolvimento das Tecnologias}

Nesta seção são apresentados os aspectos relacionados às tecnologias adotadas no desenvolvimento do sistema. O objetivo é descrever as principais decisões técnicas envolvidas na construção da solução proposta.

\section{Tecnologias Utilizadas}
\subsection{Desenvolvimento - Next.js e TypeScript}
A construção da aplicação foi fundamentada na utilização do \textit{framework} Next.js, uma ferramenta baseada em React que oferece funcionalidades como renderização no lado do servidor (SSR) e geração de sites estáticos (SSG). A adoção do Next.js permitiu a criação de uma aplicação web moderna, performática e otimizada para SEO. O uso de TypeScript trouxe tipagem estática ao projeto, aumentando a segurança do código, facilitando a detecção de erros durante o desenvolvimento e melhorando a manutenção a longo prazo. A estrutura de rotas baseada em arquivos e os \textit{Server Actions} simplificaram a comunicação entre o cliente e o servidor, resultando em um fluxo de desenvolvimento ágil e robusto.

\subsection{Estilização e Interface - Tailwind CSS}
Para a construção da interface gráfica da aplicação, adotou-se o \textit{framework} Tailwind CSS, uma ferramenta baseada em utilitários que permite a criação de designs personalizados diretamente no HTML. A utilização do Tailwind CSS proporcionou um desenvolvimento visual extremamente rápido e consistente, eliminando a necessidade de escrever CSS tradicional complexo. O sistema de design integrado facilitou a criação de uma interface responsiva, moderna e altamente customizável, garantindo uma excelente experiência de usuário em diferentes dispositivos, desde smartphones até computadores de mesa.

\subsection{Banco de Dados - MongoDB}
Para a camada de persistência de dados, foi adotado o MongoDB, um banco de dados NoSQL orientado a documentos. Sua escolha justifica-se pela flexibilidade no armazenamento de dados semiestruturados, como as questões do ENEM e os registros de desempenho dos alunos. A integração com o Next.js foi realizada de forma eficiente, permitindo operações rápidas de leitura e escrita. O modelo de dados em formato JSON (BSON) do MongoDB alinha-se perfeitamente com a estrutura de objetos do JavaScript/TypeScript utilizada no projeto, simplificando a camada de acesso a dados e garantindo escalabilidade para o crescimento da base de questões e usuários.

\subsection{Inteligência Artificial - API Groq com Modelo Llama 3}
A integração com Inteligência Artificial foi realizada por meio da API Groq \cite{groq}, utilizada para acesso a modelos de linguagem de grande escala (LLMs). A comunicação com a API foi implementada no servidor utilizando a biblioteca oficial da Groq para Node.js, enviando as questões e recebendo os feedbacks gerados pelo modelo \textit{Llama 3.3 70B Versatile}, desenvolvido pela Meta AI. A infraestrutura de aceleração da Groq permitiu tempos de inferência extremamente baixos, proporcionando feedbacks quase instantâneos para os estudantes.

% ------------------------------------------------
% DISCUSSÃO TÉCNICA DA IMPLEMENTAÇÃO
% ------------------------------------------------

\chapter{Discussão Técnica da Implementação}

O desenvolvimento resultou em um protótipo funcional da plataforma NextEducation, construído com o framework Next.js (utilizando a arquitetura App Router) e estilizado nativamente com Tailwind CSS, focando em uma interface moderna (Dark Mode) que prioriza a concentração do estudante. O sistema é composto por módulos integrados que guiam o aluno desde o primeiro acesso até a análise de desempenho:

\begin{itemize}
    \item \textbf{Painel Geral (Dashboard)}: A interface principal (Figura \ref{fig:dashboard}) atua como um \textit{hub} central. Ela apresenta métricas em tempo real, como o volume de questões indexadas no MongoDB a partir do dataset do ENEM, \textit{baseline} de proficiência e atalhos rápidos para iniciar os estudos ou realizar simulados. A comunicação com o banco de dados ocorre de forma assíncrona utilizando rotas internas (API Routes).
    
    \item \textbf{Módulos de Estudos e Simulados}: Áreas dedicadas onde o aluno pode resolver questões filtradas por área de conhecimento, especificamente Ciências Humanas. O frontend captura as respostas do usuário e as envia para o servidor (ex: \texttt{/api/questions}), que gerencia a validação do gabarito. (Figura \ref{fig:simulado})
    
    \item \textbf{Tutor IA Integrado}: O núcleo inteligente da plataforma. Ao errar uma questão ou solicitar explicação, o sistema aciona as rotas na pasta \texttt{/api/ai} que se comunicam com a API da Groq. O modelo recebe o contexto da questão, a alternativa escolhida pelo aluno e o gabarito oficial, retornando um feedback pedagógico detalhado. (Figura \ref{fig:tutor})
\end{itemize}

\begin{figure}[H]
    \centering
    \includegraphics[width=0.9\textwidth]{Dash.png}
    \caption{Print da tela inicial (Dashboard) da plataforma NextEducation.}
    \label{fig:dashboard}
\end{figure}

\begin{figure}[H]
    \centering
    \includegraphics[width=0.9\textwidth]{Simulado.png}
    \caption{Print da interface do módulo de simulados e atividades.}
    \label{fig:simulado}
\end{figure}

\begin{figure}[H]
    \centering
    \includegraphics[width=0.9\textwidth]{Tutor.png}
    \caption{Print da interação com o Tutor IA e geração de feedback didático.}
    \label{fig:tutor}
\end{figure}

Para avaliar a integração com a API Groq \cite{groq}, foi executado um conjunto de testes preliminares com questões de Ciências Humanas extraídas do dataset ENEM-Benchmark \cite{maritaca}. Cada questão foi submetida ao sistema juntamente com a alternativa selecionada pelo usuário e o gabarito correto. O modelo Llama 3.3 foi instruído via engenharia de \textit{prompt} a atuar como um professor particular, gerando um feedback didático explicando por que a alternativa selecionada estava incorreta e justificando o gabarito oficial.

\section{Resultados Parciais}

Os testes do protótipo, abrangendo a amostra de 20 questões de Ciências Humanas avaliadas consensualmente pelos autores, foram analisados em duas dimensões principais: estabilidade técnica e qualidade pedagógica. A Tabela \ref{tab:resultados} resume as médias observadas nos critérios de validação da IA.

\begin{table}[H]
    \centering
    \caption{Resultados Preliminares da Avaliação Pedagógica da IA (Escala 1 a 5)}
    \begin{tabular}{|l|c|}
        \hline
        \textbf{Critério de Avaliação} & \textbf{Média Obtida} \\
        \hline
        Correção Conceitual & 4,5 \\
        Aderência ao Gabarito & 4,8 \\
        Clareza Didática & 4,2 \\
        Contextualização em Ciências Humanas & 4,6 \\
        \hline
    \end{tabular}
    \label{tab:resultados}
\end{table}

Em termos de \textbf{estabilidade técnica}, o uso do framework Next.js em conjunto com a API Groq demonstrou resultados promissores. O tempo médio de resposta para a geração de feedbacks complexos foi mantido na casa dos milissegundos a poucos segundos, caracterizando uma latência mínima, fator essencial para manter o engajamento durante o estudo ativo. A infraestrutura técnica lidou adequadamente com as requisições assíncronas à API e as operações de leitura/escrita no MongoDB.

No âmbito da \textbf{qualidade pedagógica}, os resultados parciais revelam que as explicações geradas pelo modelo são textualmente coesas, com forte aderência ao gabarito e boa contextualização histórica e sociológica. O sistema conseguiu, na maioria dos casos testados, isolar o conceito que o aluno demonstrou não compreender e fornecer uma explicação direcionada.

Entretanto, conforme indicado na pontuação mais baixa de "Clareza Didática", observou-se a necessidade de refinamentos iterativos no \textit{prompt} do sistema. Em algumas interações, o modelo tendeu a fornecer explicações excessivamente genéricas ou verborrágicas, o que pode dispersar a atenção do aluno. Além disso, embora a taxa de alucinação tenha sido baixa nas respostas de Ciências Humanas, ela não é nula. Assim, reforça-se que o sistema deve ser compreendido como uma ferramenta complementar de avaliação formativa e apoio ao estudo, e não como um substituto absoluto para a mediação docente.

% ------------------------------------------------
% CONCLUSÃO
% ------------------------------------------------

\chapter{Conclusão}

Este trabalho apresentou o desenvolvimento da plataforma NextEducation, um sistema web de apoio ao estudo de Ciências Humanas para o ENEM, integrado a modelos de linguagem por meio da API Groq. A plataforma reúne módulos de autenticação, atividades, mini provas, dashboard de desempenho e tutor IA, buscando transformar a correção de questões objetivas em uma experiência de avaliação formativa.

A principal contribuição do trabalho está na integração entre questões do ENEM, análise de desempenho e geração automática de feedback didático. Essa combinação permite que o estudante não apenas identifique se acertou ou errou uma questão, mas receba uma explicação contextualizada sobre seu desempenho.

A implementação demonstrou viabilidade técnica inicial, especialmente na integração entre a plataforma web e a API Groq. No entanto, a avaliação ainda precisa ser ampliada por meio de testes sistemáticos com maior quantidade de questões, análise da qualidade pedagógica das respostas e participação de estudantes ou professores na avaliação do feedback gerado.
Como trabalhos futuros, recomenda-se realizar uma avaliação empírica com estudantes, comparar o feedback gerado pela IA com explicações elaboradas por professores, analisar a presença de erros conceituais ou alucinações e investigar o impacto da plataforma na aprendizagem ao longo do tempo.


% ------------------------------------------------
% REFERÊNCIAS
% ------------------------------------------------

\postextual

\begin{thebibliography}{99}

\bibitem{ebbinghaus}
EBBINGHAUS, H. \textbf{Memory: A Contribution to Experimental Psychology}. New York: Teachers College, Columbia University, 1913.

\bibitem{cepeda}
CEPEDA, N. J. et al. \textbf{Distributed practice in verbal recall tasks: A review and quantitative synthesis}. Psychological Bulletin, v. 132, n.3, p.354–380, 2006.

\bibitem{kang}
KANG, S. H. K. \textbf{Spaced repetition promotes efficient and effective learning: Policy implications for instruction}. Policy Insights from the Behavioral and Brain Sciences, v.3, n.1, p.12–19, 2016.

\bibitem{kasneci2023chatgpt}
KASNECI, E. et al. \textbf{ChatGPT for Good? On Opportunities and Challenges of Large Language Models for Education}. Learning and Individual Differences, v. 103, artigo 102274, p.1--9, 2023.

\bibitem{maritaca}
MARITACA AI. \textbf{ENEM-Benchmark: A Multi-modal Dataset for Evaluating LLMs}. Versão 1.0. Hugging Face, 2024. Disponível em: <\url{https://huggingface.co/datasets/maritaca-ai/enem}>. Acesso em: 25 mar. 2026.

\bibitem{zhao2025lector}
ZHAO, H. et al. \textbf{LECTOR: LLM-Enhanced Concept-based Test-Oriented Repetition for Adaptive Spaced Learning}. \textit{arXiv preprint arXiv:2508.03275}, 2025. Disponível em: <\url{https://arxiv.org/abs/2508.03275}>. Acesso em: 9 mar. 2026.

\bibitem{damasceno2021stuart}
DAMASCENO, Adson Roberto. \textbf{STUART: Um Sistema de Tutoria Inteligente Artificial para aumentar a escalabilidade dos cursos a distância}. 2021. Dissertação (Mestrado em Ciência da Computação) -- Universidade Estadual do Ceará (UECE), Fortaleza, 2021.

\bibitem{seabra2023memorizapp}
SEABRA, Igor. \textbf{MemorizApp: um sistema de aprendizado baseado em repetição espaçada}. 2023. Trabalho de Conclusão de Curso (Graduação em Ciência da Computação) -- Universidade Federal de Campina Grande (UFCG), Campina Grande, 2023.

\bibitem{groq}
GROQ. API Reference. GroqDocs, 2026. Disponível em: <\url{https://console.groq.com/docs/api-reference}>. Acesso em: 27 mai. 2026.

\bibitem{nextjs}
NEXT.JS. \textbf{Next.js Documentation}. Vercel, 2026. Disponível em: <\url{https://nextjs.org/docs}>. Acesso em: 6 mai. 2026

\bibitem{typescript}
TYPESCRIPT. \textbf{TypeScript Documentation}. Microsoft, 2026. Disponível em: <\url{https://www.typescriptlang.org/docs/}>. Acesso em: 6 mai. 2026

\bibitem{tailwindcss}
TAILWIND CSS. \textbf{Tailwind CSS Documentation}. Tailwind Labs, 2026. Disponível em: <\url{https://tailwindcss.com/docs}>. Acesso em: 6 mai. 2026.

\bibitem{mongodb}
MONGODB. \textbf{MongoDB Documentation}. MongoDB Inc., 2026. Disponível em: <\url{https://www.mongodb.com/docs/}>. Acesso em: 6 mai. 2026.

\bibitem{blackwilliam}
BLACK, P.; WILIAM, D. \textbf{Assessment and classroom learning}. Assessment in Education: Principles, Policy and Practice, v.5, n.1, p.7–74, 1998.

\bibitem{sadler}
SADLER, D. R. \textbf{Formative assessment and the design of instructional systems}. Instructional Science, v.18, p. 119–144, 1989.

\end{thebibliography}

\end{document}